\section{Marco teorico}

% TODO: desarrollar cada subseccion. Las formulas ya estan planteadas y
% coinciden con las docstrings del codigo fuente; el texto debe explicarlas.

\subsection{La imagen como objeto algebraico}

Una imagen digital en color es un tensor $I \in \mathbb{R}^{H \times W \times 3}$.
% TODO: desarrollar.

\subsection{Transformaciones lineales de color}

La conversion de RGB a YCbCr, en la recomendacion ITU-R BT.601, es una
transformacion afin cuya parte lineal es una matriz de $3 \times 3$:

\begin{equation}
\begin{bmatrix} Y \\ C_b \\ C_r \end{bmatrix}
=
\begin{bmatrix}
 0.2990 &  0.5870 &  0.1140 \\
-0.1687 & -0.3313 &  0.5000 \\
 0.5000 & -0.4187 & -0.0813
\end{bmatrix}
\begin{bmatrix} R \\ G \\ B \end{bmatrix}
+
\begin{bmatrix} 0 \\ 128 \\ 128 \end{bmatrix}.
\end{equation}

La conversion de RGB a HSV, en cambio, \emph{no} es lineal: el tono y la
saturacion se definen a partir de $\max(R,G,B)$, $\min(R,G,B)$ y cocientes
entre esas cantidades, operaciones que no admiten representacion matricial.
Esa es la razon por la que la segmentacion de este trabajo opera en YCbCr.
% TODO: ampliar.

\subsection{Momentos y matriz de inercia}

Con los momentos centrales de segundo orden de la mascara binaria se construye

\begin{equation}
M = \frac{1}{\mu_{00}}
\begin{bmatrix} \mu_{20} & \mu_{11} \\ \mu_{11} & \mu_{02} \end{bmatrix},
\end{equation}

matriz simetrica y semidefinida positiva cuyos autovectores son las direcciones
principales del objeto y cuyos autovalores miden la dispersion a lo largo de
esas direcciones.
% TODO: desarrollar.

\subsection{Descomposicion en valores singulares y eigen-objetos}

Sea $A \in \mathbb{R}^{m \times d}$ la matriz de imagenes centradas. Su
descomposicion en valores singulares $A = U \Sigma V^{T}$ entrega en las
columnas de $V$ los autovectores de la matriz de covarianza
$C = A^{T}A/(m-1)$. Cuando $m \ll d$ se aplica la identidad de
\textcite{turk1991}: si $A A^{T} u = \mu\, u$, entonces

\begin{equation}
A^{T}A \,(A^{T}u) = \mu \,(A^{T}u),
\end{equation}

de modo que $v = A^{T}u / \lVert A^{T}u \rVert$ es autovector de la matriz
grande sin necesidad de construirla.
% TODO: desarrollar.

\subsection{Minimos cuadrados regularizados}

El clasificador resuelve

\begin{equation}
\min_{W} \; \lVert XW - Y \rVert_{F}^{2} + \lambda \lVert W \rVert_{F}^{2},
\qquad
W = (X^{T}X + \lambda I)^{-1} X^{T} Y ,
\end{equation}

que con $\lambda = 0$ coincide con $W = X^{+}Y$, donde $X^{+}$ es la
pseudoinversa de Moore-Penrose \parencite{boyd2018}.
% TODO: desarrollar, incluir la interpretacion geometrica de la proyeccion
% sobre el espacio columna de X y la frontera de decision como hiperplano.

\subsection{Numero de condicion y regularizacion}

\begin{equation}
\operatorname{cond}_2(X^{T}X + \lambda I)
= \frac{\sigma_{\max} + \lambda}{\sigma_{\min} + \lambda}
\end{equation}

decrece monotonamente con $\lambda$: la regularizacion desplaza todo el
espectro y estabiliza la solucion a costa de sesgo.
% TODO: desarrollar.

\subsection{Mapeo de conceptos}

\begin{table}[htbp]
\caption{Conceptos de algebra lineal y su punto de aplicacion en el sistema}
\label{tab:mapeo-conceptos}
\begin{tabular}{ll}
\toprule
Concepto & Donde se aplica \\
\midrule
Matrices y tensores & Imagen como $I \in \mathbb{R}^{H \times W \times 3}$ \\
Producto matriz-vector & Conversion RGB a YCbCr y RGB a gris \\
Autovalores y autovectores & Matriz de inercia $2\times2$, ejes principales \\
Matriz de covarianza & Estandarizacion y analisis de componentes principales \\
Descomposicion en valores singulares & Eigen-objetos, reduccion de 4096 a $k$ \\
Relacion entre $A^{T}A$ y $AA^{T}$ & Truco de Turk-Pentland para $m \ll n$ \\
Espacios y subespacios vectoriales & Espacio de caracteristicas, subespacio propio \\
Proyeccion ortogonal & Representacion en el subespacio de eigen-objetos \\
Minimos cuadrados & Entrenamiento del clasificador \\
Pseudoinversa de Moore-Penrose & Solucion cerrada $W = X^{+}Y$ \\
Rango e independencia lineal & Deteccion de caracteristicas redundantes \\
Numero de condicion & Diagnostico de $X^{T}X$ mal condicionada \\
Regularizacion de Tikhonov & Estabilizacion mediante $+\lambda I$ \\
Norma euclidiana & Baseline k-NN, distancia en el espacio propio \\
Hiperplanos separadores & Frontera de decision del clasificador lineal \\
\bottomrule
\end{tabular}
\note{Elaboracion propia.}
\end{table}

